\chapter{Methodology}
\label{ch:methodology}

This chapter describes the complete technical approach taken to build, train, and evaluate
the vulnerability detection system.
It is organised into five parts: model architecture, dataset construction, training procedure,
data augmentation strategies, and the evaluation framework.

\section{Model Architecture}

The base model is UniXcoder (\texttt{microsoft/unixcoder-base}), a 125-million-parameter
encoder-only transformer pre-trained jointly on code and natural language across multiple
programming languages~\cite{guo2022unixcoder}.
UniXcoder extends the RoBERTa architecture with cross-modal pre-training objectives including
masked language modelling over code tokens, natural language descriptions, and abstract
syntax tree linearisations.

For the vulnerability detection task, a lightweight classification head is attached to the
UniXcoder encoder.
The encoder produces a sequence of contextual token embeddings; the embedding corresponding
to the special \texttt{[CLS]} token, positioned at the start of every input, is used as
the fixed-length representation of the entire function.
This representation is passed through a dropout layer (dropout rate $p = 0.1$) before being
projected to a single scalar by a fully-connected linear layer ($d_{\text{hidden}} \to 1$).
A sigmoid activation converts this scalar to a probability in $[0, 1]$, interpreted as
the probability that the input function is vulnerable.
The classification threshold is set to 0.5 by default.

Classifier weights are initialised from a Gaussian distribution with mean zero and
standard deviation 0.02.
All other parameters are initialised from the UniXcoder checkpoint.

\subsection{Label Smoothing}

Standard binary cross-entropy training assigns hard target labels of exactly 0 (safe) or
1 (vulnerable).
To prevent overconfidence, label smoothing with $\varepsilon = 0.1$ is applied during training.
The hard targets are replaced by soft targets:
\begin{equation}
  \tilde{y} = y \cdot (1 - \varepsilon) + \varepsilon \cdot 0.5
\end{equation}
A vulnerable example ($y = 1$) becomes target $\tilde{y} = 0.95$;
a safe example ($y = 0$) becomes $\tilde{y} = 0.05$.
This regularisation prevents the model from being penalised for moderate confidence scores,
which is desirable in security applications where calibrated uncertainty is valuable.

\section{Dataset Construction}

\subsection{Vulnerability Classes}

The dataset covers 17 CWE types across C and Python.
C-specific classes include: buffer overflow (CWE-120), use-after-free (CWE-416),
double free (CWE-415), integer overflow (CWE-190), null pointer dereference (CWE-476),
off-by-one error (CWE-193), and uninitialized variable read (CWE-457).
Python-specific classes include: insecure deserialisation (CWE-502),
server-side request forgery (CWE-918), and eval injection (CWE-95).
Eight classes are shared across both languages: command injection (CWE-78),
TOCTOU race condition (CWE-367), hardcoded credentials (CWE-798), path traversal (CWE-22),
insecure file permissions (CWE-732), weak cryptography (CWE-327), SQL injection (CWE-89),
and format string or template injection (CWE-134).

These 17 classes were chosen to balance three criteria.
First, industrial prevalence: all selected classes appear in the CWE Top 25 Most Dangerous
Software Weaknesses or are frequently represented in real-world CVE disclosures.
Second, structural diversity: the set deliberately spans locally detectable patterns
(CWE-798 hardcoded credentials, CWE-327 weak cryptography) and multi-statement patterns
(CWE-416 use-after-free, CWE-367 TOCTOU), so the evaluation distinguishes what
sequence models handle well from where they fail.
Third, cross-language coverage: eight classes are shared between C and Python, enabling
the ablation study to measure transfer quality directly rather than inferring it from
aggregate scores.
This selection is not exhaustive --- the CWE landscape contains over 900 entries --- but
it is representative of the classes most likely to appear in production codebases and most
informative for evaluating a sequence-based detector.

Table~\ref{tab:cwe_dataset} summarises all 17 classes by language scope and pattern type.

\begin{table}[H]
  \centering
  \small
  \caption{The 17 CWE vulnerability classes used in the dataset, grouped by language scope.}
  \label{tab:cwe_dataset}
  \begin{tabular}{llll}
    \toprule
    Language Scope & CWE ID & Vulnerability Name & Pattern Type \\
    \midrule
    \multirow{7}{*}{C only}
      & CWE-120 & Buffer Overflow              & Memory Safety \\
      & CWE-416 & Use-After-Free               & Memory Safety \\
      & CWE-415 & Double Free                  & Memory Safety \\
      & CWE-190 & Integer Overflow             & Memory Safety \\
      & CWE-476 & Null Pointer Dereference     & Memory Safety \\
      & CWE-193 & Off-by-One Error             & Memory Safety \\
      & CWE-457 & Uninitialized Variable Read  & Memory Safety \\
    \midrule
    \multirow{3}{*}{Python only}
      & CWE-502 & Insecure Deserialisation     & Deserialisation \\
      & CWE-918 & Server-Side Request Forgery  & Web / Runtime \\
      & CWE-95  & Eval Injection               & Injection \\
    \midrule
    \multirow{8}{*}{Shared (C \& Python)}
      & CWE-78  & Command Injection            & Injection \\
      & CWE-367 & TOCTOU Race Condition        & Race Condition \\
      & CWE-798 & Hardcoded Credentials        & Access Control \\
      & CWE-22  & Path Traversal               & Access Control \\
      & CWE-732 & Insecure File Permissions    & Access Control \\
      & CWE-327 & Weak Cryptography            & Cryptographic \\
      & CWE-89  & SQL Injection                & Injection \\
      & CWE-134 & Format String / Template Inj.& Injection \\
    \bottomrule
  \end{tabular}
\end{table}

Each vulnerable example is a short self-contained function illustrating the pattern precisely.
Each safe counterpart uses the same structural template but applies the correct mitigation:
parameterised queries instead of string concatenation for CWE-89, environment-variable lookup
instead of a hardcoded string for CWE-798, \texttt{hashlib.sha256} instead of
\texttt{hashlib.md5} for CWE-327.
Keeping the mitigated code structurally close to the vulnerable code is deliberate:
a model that genuinely understands vulnerability patterns should detect the specific dangerous
operation, not associate risk with superficial features such as function length or identifier choice.

\subsection{Contextual Wrapping}

Approximately half of all training examples are wrapped in surrounding caller context using a
\texttt{wrap\_in\_context()} function applied at dataset generation time.
The wrapper adds a plausible outer function that calls the vulnerable or safe inner function,
simulating the way real code embeds helper routines inside larger call chains.
The wrapping probability is 0.5, so the model sees isolated functions and context-embedded
snippets in roughly equal proportion.

\subsection{Dataset Splits}

The generated synthetic data is split into training, validation, and held-out test sets.
The held-out validation set (\texttt{multi\_lang\_valid.jsonl}) contains 1{,}230 samples:
660 C examples and 570 Python examples, with a roughly balanced label distribution.
This split is used for all quantitative evaluation in Experiments 1--4 and is never used
during training or hyperparameter selection.
The real-world evaluation in Experiment~3 uses \texttt{sample\_test.jsonl} from
DiverseVul~\cite{ni2023diversevul}, which contains 200 real-world C functions and is
never exposed to training.

\section{Training Procedure}

Three model variants are trained to support the cross-language ablation study:
\begin{itemize}
  \item \textbf{C-only model}: trained exclusively on C examples from the synthetic dataset.
  \item \textbf{Python-only model}: trained exclusively on Python examples from the synthetic dataset.
  \item \textbf{Joint model}: trained on the combined C and Python dataset, interleaved and shuffled.
\end{itemize}

All three variants use identical hyperparameters.
The optimiser is AdamW with initial learning rate $2 \times 10^{-5}$ and weight decay
$\lambda = 0.01$.
Bias terms and layer normalisation parameters are excluded from weight decay.
A linear learning rate warm-up covering 10\% of total training steps is followed by linear
decay to zero.
Training runs for three epochs with a per-GPU batch size of 8.
The maximum tokenised sequence length is 400 tokens; functions exceeding this limit are
truncated to the first 400 tokens.
The checkpoint achieving the lowest validation loss is saved as the best model.

\section{Data Augmentation Strategies}

\subsection{Hard Negative Mining}

Hard negatives are safe examples that are deliberately designed to resemble vulnerable ones
in surface form while being semantically safe.
The objective is to force the model to learn the \emph{discriminating} feature rather than
the \emph{correlated} feature.

Six hard negative examples are added to the Python safe set: functions using
\texttt{hashlib.sha256}, \texttt{hashlib.sha512}, and \texttt{hashlib.pbkdf2\_hmac} in the
same structural positions where vulnerable examples use \texttt{hashlib.md5}.
These safe examples use identical code structure to the vulnerable ones, differing only in
the algorithm name, creating a minimal contrastive pair.

Four safe examples are added to the C safe set: functions calling \texttt{system()} with
fully hardcoded string literals as arguments, with no user-controlled data in the call chain.
Additional safe examples are added for OpenSSL SHA-256 and SHA-512 usage.
In total, seventeen new safe examples are added across both languages.

\subsection{Real-World Data Augmentation}

To reduce the synthetic-to-real generalisation gap, 1{,}000 function-level samples from
the DiverseVul training split (\texttt{sample\_train.jsonl}) are incorporated into the
joint model's training data.
Each real-world record is normalised to the training schema: the function body (\texttt{func}),
binary label (\texttt{target}), a unique identifier of the form \texttt{real\_<idx>},
language tag \texttt{c}, and available CWE annotations.

The 1{,}000 real samples are merged with the synthetic training set, shuffled, and used for
a single fine-tuning run.
The DiverseVul test split (\texttt{sample\_test.jsonl}, 200 samples) is kept completely
separate and used only for Experiment~3 evaluation; no sample from the test split is
exposed to training at any point.

\section{Evaluation Framework}

Four quantitative experiments and one qualitative test suite are used to evaluate the model.

\textbf{Experiment 1 --- Baseline performance} evaluates the joint model on the held-out
synthetic validation set.
Primary metrics are F1-score, precision, recall, and AUC-ROC.
Per-language and per-CWE breakdowns are computed.

\textbf{Experiment 2 --- Cross-language ablation} tests C-only, Python-only, and Joint
models against C-only, Python-only, and Combined held-out test sets.
Each combination produces an F1-score, forming a 3$\times$3 matrix.

\textbf{Experiment 3 --- Real-world generalisation} evaluates the joint model on the
200-sample DiverseVul test split.
Synthetic and real-world F1 are compared to measure the generalisation gap.
This experiment is run both before and after the real-world data augmentation.

\textbf{Experiment 4 --- Threshold sweep} evaluates how varying the decision threshold from
0.1 to 0.9 affects precision, recall, and F1 on the synthetic validation set.

\textbf{Semantic understanding tests} probe whether the model is pattern-matching or
performing genuine semantic reasoning, using four test types:

\begin{enumerate}
  \item \emph{Variable renaming}: functions with renamed identifiers but identical vulnerability
        logic. If prediction changes after renaming, the model is sensitive to identifier names.
  \item \emph{Minimal fix}: pairs of vulnerable and fixed functions where only the dangerous
        operation is changed. A correct model should flip its prediction on the fixed version.
  \item \emph{Dead code injection}: safe functions with unreachable vulnerable blocks guarded
        by \texttt{if (0)} or \texttt{if False}. A model without control-flow reasoning will
        be misled by dangerous tokens in dead branches.
  \item \emph{Token ablation}: the single identifying keyword (e.g., \texttt{system},
        \texttt{md5}) is replaced with a neutral name. If the model still detects the
        vulnerability from context, it demonstrates understanding beyond keyword matching.
\end{enumerate}
